\section*{Chapter \ref{ch:s2:the-bond-futures-basis-trade} --- The Bond-Futures Basis Trade}

\begin{solution}{exo:s2:the-bond-futures-basis-trade:1}
$1/(0.01 + 0.01) = 50$ times; at 3\% each, $1/0.06 = 16.7$ times.
\end{solution}

\begin{solution}{exo:s2:the-bond-futures-basis-trade:2}
$0.002 \times 50 = 10\%$ a year on capital, before costs and compounding.
\end{solution}

\begin{solution}{exo:s2:the-bond-futures-basis-trade:3}
$50 \times 0.006 = 30\%$ of capital.
\end{solution}

\begin{solution}{exo:s2:the-bond-futures-basis-trade:4}
At ten times, capital covers the stressed collateral (6\% of face needs 0.6 of capital), so the book keeps its whole position and gets the basis back.
At fifty times it must sell down to 16.7 times at the stressed basis; the positions it sold never recover, and the sales cost 10 cents per 100.
\end{solution}

\begin{solution}{exo:s2:the-bond-futures-basis-trade:5}
$(659 - 554)/659 = 15.9\%$.
\end{solution}

\begin{solution}{exo:s2:the-bond-futures-basis-trade:6}
Leveraged funds hold Treasury futures for many reasons (duration views, curve trades, trend following), and the report shows futures only, not the
cash bonds or repo on the other side; a net short can be a basis trade or an outright bet.
\end{solution}

\begin{solution}{exo:s2:the-bond-futures-basis-trade:7}
The fifty-times book still sells on 66 days, since it sits at its collateral cap; its trough is deeper ($-27.7\%$) because it cuts less early, and it
keeps more of the recovery ($-6.2\%$ after instead of $-9.6\%$). Part of the locked-in loss comes from the higher collateral terms, part from running
at the cap.
\end{solution}

\begin{solution}{exo:s2:the-bond-futures-basis-trade:8}
Convergence at delivery helps only a book that is still holding at delivery. Margin calls and repo that cannot be rolled can force a sale before
then, at the basis of the day, and the loss is then real.
\end{solution}

\begin{solution}{pb:s2:the-bond-futures-basis-trade:1}
\begin{enumerate}
\item Gross basis: bond price minus futures price times conversion factor; net basis: gross basis less carry; implied repo: the financing rate at which
buying, delivering and repaying breaks even.
\item When the implied repo rate exceeds the repo rate paid.
\item Face per unit of capital, capped by one over haircut plus margin; the risk that repo cannot be refinanced on the same terms.
\item The collateral: haircut plus futures margin.
\item They want duration without using their balance sheets and buy futures.
\item \$400 to \$500 billion at the peak, over 60\% of hedge funds' Treasury exposure and 70\% of their repo borrowing.
\item Leveraged funds net short \$475 billion on 18 February 2020, a record \$1\,184 billion on 12 November 2024.
\item They hold Treasuries for futures longs and fund them in repo.
\item Wider bid--ask spreads, repo rates jumping, arbitrage spreads diverging, higher futures margins.
\item Cut short futures from \$659 to \$554 billion and sold \$91 to \$105 billion of Treasuries.
\item At least \$500 billion of Treasury purchases on 15 March, then purchases in the amounts needed on 23 March.
\item No: the Treasuries the trade held kept trading at a premium before 17 March.
\item \textbf{Named result.} The book earns 2.1\% a year at ten times and 10.3\% at fifty; in the planted funding stress the fifty-times book is forced
to sell and ends 9.6\% below its pre-stress capital after the basis recovers, the ten-times book 0.2\%.
\item Marks $-6.6\%$, forced-sale costs $-3.3\%$, dearer repo $-0.6\%$, carry $+0.8\%$.
\item Term repo keeps the funding in place through the stress, so only the futures margin can force sales.
\item Futures positions of leveraged funds, repo volumes, the futures-implied versus cash yield gap, collateral terms.
\item So that stressed collateral terms and a basis move leave capital above the requirement.
\item Stress collateral terms, include forced-sale costs, use the repo actually available.
\item The basis-trade crowding signal.
\item Balance sheet: it holds Treasuries for those who want futures.
\end{enumerate}
\end{solution}

\begin{solution}{iq:s2:the-bond-futures-basis-trade:1}
Futures trade rich to cash when demand for them is high; a trader buys the cheapest-to-deliver bond, finances it in repo, sells the future and earns
the implied repo rate less the repo rate, at high leverage because both legs are collateralised.
\end{solution}

\begin{solution}{iq:s2:the-bond-futures-basis-trade:2}
A joint rise in haircut and margin, a basis move against the book and a repo spike, applied to the whole position at once; set leverage so that
capital still covers collateral under the stressed terms, and hold term funding.
\end{solution}

\begin{solution}{iq:s2:the-bond-futures-basis-trade:3}
From the bond's full price today, the invoice price at delivery (futures times conversion factor plus accrued) and the coupons in between; a gap above
repo means the futures are rich, below means they are cheap.
\end{solution}

\begin{solution}{iq:s2:the-bond-futures-basis-trade:4}
Post more capital if available; otherwise reduce the position, ideally in the less liquid legs first, and seek term repo; never let the margin call
go unmet.
\end{solution}

\begin{solution}{iq:s2:the-bond-futures-basis-trade:5}
Prices and repo rates each morning; collateral requirements from each lender and the exchange; variation margin; the capital check against stressed
terms; roll schedules; with alerts when capital approaches the requirement.
\end{solution}

\begin{solution}{iq:s2:the-bond-futures-basis-trade:6}
After a move $x$ capital is $C - Nx$ and the requirement $Nk$, so sales start when $C - Nx < Nk$. With $N = LC$ the largest move without selling is
$x < 1/L - k$: at $L = 50$ and $k = 0.02$ it is zero, which is why a book at its cap sells on ordinary days.
\end{solution}
